\begin{exercisechunk}
\item \textbf{A large penalty and two-stage fitting.}
\label[exercise]{ex:l9-two-stage-fitting}
\exerciseprofile{3}{2}{3}{\exproof\quad\excalculation\quad\exinterpretation}{%
Quantify when a weighted interpolator approximates fitting the first block and then its residual.}
Fix $U\in\R^{n\times m}$ of column rank $m$,
$V\in\R^{n\times D}$ of row rank $n$, and $\mathbf Y\in\R^n$.
For $\lambda>0$, let $(\widehat\beta_\lambda,\widehat b_\lambda)$ minimize
$\|\beta\|^2+\lambda\|b\|^2$ subject to $U\beta+Vb=\mathbf Y$.
Set $H=D^{-1}VV^\top$, $A=U^\top H^{-1}U$, and define the two-stage fit by
\[
 \widetilde\beta
 =\arg\min_\beta(\mathbf Y-U\beta)^\top H^{-1}(\mathbf Y-U\beta),
 \qquad
 \widetilde b=V^\top(VV^\top)^{-1}(\mathbf Y-U\widetilde\beta).
\]
\begin{enumerate}[(a),beginpenalty=10000]
\item Prove that
\[
 \widetilde\beta=A^{-1}U^\top H^{-1}\mathbf Y,
 \qquad
 \widehat\beta_\lambda-\widetilde\beta
 =-\bigl(I_m+(\lambda/D)A\bigr)^{-1}\widetilde\beta.
\]
If $A\succeq aI_m$ with $a>0$, deduce
\[
 \|\widehat\beta_\lambda-\widetilde\beta\|
 \le\frac{\|\widetilde\beta\|}{1+(\lambda/D)a}.
\]

\item Put $d_\lambda=\widehat\beta_\lambda-\widetilde\beta$.
Prove that
\[
 \widehat b_\lambda-\widetilde b
 =-V^\top(VV^\top)^{-1}Ud_\lambda,
 \qquad
 \|\widehat b_\lambda-\widetilde b\|^2=D^{-1}d_\lambda^\top A d_\lambda.
\]
For the lecture's choice $\lambda=D$, show that on $\cG$,
\[
 \|\widehat\beta_D-\widetilde\beta\|
 \le\frac{\|\widetilde\beta\|}{1+n/(2C_H)},\qquad
 \|\widehat b_D-\widetilde b\|^2
 \le\frac{2C_H}{Dn}\|\widetilde\beta\|^2.
\]
Distinguish the limit $\lambda\to\infty$ with $U,V,\mathbf Y$ fixed from
the lecture's bound, where $D$, $V$, and $n$ may vary together.

\item Show that $\widetilde\beta$ is the ordinary least-squares solution
when $H$ is a scalar multiple of $I_n$. For
\[
 U=\begin{pmatrix}1\\1\end{pmatrix},\quad
 V=\begin{pmatrix}\sqrt2&0\\0&2\sqrt2\end{pmatrix},\quad
 \mathbf Y=\begin{pmatrix}0\\1\end{pmatrix},\quad D=2,
\]
calculate $\widetilde\beta$ and the ordinary least-squares solution.
Explain why a large second-block penalty alone does not imply convergence
to ordinary least squares.
\end{enumerate}
\begin{hints}
Eliminate $b$ as in \cref{ex:l9-minimum-norm}.
For the last bound in (b), diagonalize $A$ and use
$t/(1+t)^2\le1/t$ for $t>0$.
\end{hints}
\end{exercisechunk}
